% Generated draft: requires mathematical and source review.
\documentclass[11pt]{article}
\usepackage[margin=25mm]{geometry}
\usepackage{fontspec}
\setmainfont{Latin Modern Roman}
\usepackage{amsmath,amssymb,mathtools,unicode-math}
\setmathfont{Latin Modern Math}
\usepackage{enumitem,needspace,xurl,hyperref}
\setlength{\parindent}{0pt}
\setlength{\parskip}{5pt}
\setlength{\emergencystretch}{3em}
\newcommand{\N}{\mathbb{N}}
\newcommand{\Z}{\mathbb{Z}}
\newcommand{\Q}{\mathbb{Q}}
\newcommand{\R}{\mathbb{R}}
\newcommand{\C}{\mathbb{C}}
\newcommand{\F}{\mathbb{F}}
\newcommand{\A}{\mathbb{A}}
\newcommand{\PP}{\mathbb{P}}
\newcommand{\E}{\mathbb{E}}
\newcommand{\eps}{\varepsilon}
\newcommand{\dd}{\,\mathrm{d}}
\begin{document}
\section*{UnsolvedMath ID: 8800071}
\textbf{Source problem code:} catalogue entry
\subsection{Definitions and mathematical statement}
A pseudorandom-function family consists of an efficient randomized key generator and efficient evaluation algorithms \(F_k:D_\lambda\to R_\lambda\), indexed by security parameter \(\lambda\). For every probabilistic polynomial-time oracle distinguisher \(A\), require
\[
 \left|\Pr[A^{F_k}(1^\lambda)=1]-\Pr[A^{\mathcal R}(1^\lambda)=1]\right|
\]
to be negligible in \(\lambda\), where \(k\) is generated honestly and \(\mathcal R\) is a uniformly random function with the same finite domain and range. The oracle may be queried adaptively. This is the standard classical PRF interpretation; the source does not specify a quantum-query or other stronger model.

An ideal lattice is a lattice associated with an ideal of an algebraic-number ring under a specified coefficient or Minkowski embedding. ``Ideal-lattice problems'' can include different geometric worst-case approximation problems and average-case ring problems, whose parameters and reductions must be stated.

Construct an efficient PRF with a security proof based on the intended ideal-lattice hardness assumption. The excerpt does not specify a ring family, ideal problem, approximation factor, modulus/error distribution, security reduction, or concrete efficiency target. It does not demand key homomorphism or security based only on ideal-SVP; those are not added as requirements.

The summary reports ring/ideal-LWE-related PRF progress and cites Boneh--Lewi--Montgomery--Raghunathan (2013) for key-homomorphic work, while distinguishing direct ideal-SVP-only constructions as less standard. No exact theorem or reduction parameters are supplied. A ring-LWE security theorem and a direct worst-case ideal-SVP assumption must not be identified without their reductions. The 2006 request and later partial-progress assessment are retained as source reports pending current primary-literature review.
\subsection{Short English statement}
Build efficiently evaluable functions indistinguishable from random functions, with security based on a specified ideal-lattice problem. The source leaves the hardness parameters and reduction model open; later ring-LWE-based work is reported as progress, not as an identical direct ideal-SVP theorem.
\subsection{Sources}
\begin{itemize}
\item Local registry: \url{data/unsolvedmath/problems.json}, permanent UnsolvedMath ID \(8800071\), source code catalogue entry; supplied statement, background, and research summary, updated 2026-08-06. Attributed to Daniele Micciancio.
\item IPAM, \emph{Workshop on Number Theory and Cryptography---Open Problems} (2006), attachment 6240, PDF page 38: \url{https://www.ipam.ucla.edu/programs/workshops/workshop-i-number-theory-and-cryptography-open-problems/?tab=schedule}. Supplied PDF/legacy-PPT visual audit accessed 2026-07-29; it reports 19 inspected surviving attachments.
\item Boneh--Lewi--Montgomery--Raghunathan (2013), key-homomorphic PRF work cited by the summary, without full publication details or exact lattice assumptions.
\end{itemize}
\end{document}
